\newtheorem*{Lemma: only multiples of |g| contribute}{Lemma~\ref{Lemma: only multiples of |g| contribute}}
\begin{Lemma: only multiples of |g| contribute}
For a complex reflection group $W$, and a \textbf{regular} element $g\in W$, the total contribution in the sum 
\[
\sum_{\chi\in\widehat{W}}\chi(1)\cdot \chi(g^{-1})\cdot e^{-t\cdot c_{\chi}},
\]
of those characters $\chi\in\widehat{W}$ for which $c_{\chi}$ is not a multiple of $|g|$ is $0$.
\end{Lemma: only multiples of |g| contribute}

\begin{proof}
We consider the partition of the set of irreducible characters $\chi\in\widehat{W}$ into orbits under the action of $\Psi$. After Prop.~\ref{Prop: properties of Psi for Cox numbers} all characters in such an orbit have the same Coxeter number. We will show that if this Coxeter number is not a multiple of $|g|$, then the total contribution of the characters of the orbit is $0$. 

If $\chi(g)=0$ for some irreducible character $\chi$, then after Prop.~\ref{Prop: The key technical lemma}, all characters in the $\Psi$-orbit of $\chi$ evaluate $g$ to $0$.  We now deal with the remaining $\Psi$-orbits of characters $\chi$ for which $\chi(g)\neq 0$.

Consider a character $\chi$ in such an orbit and let $k$ be the smallest number such that $\Psi^k(\chi)=\chi$. Assume further that $g$ is $\zeta$-regular for an eigenvalue $\zeta=e^{2\pi il/d}$ with $(l,d)=1$ (i.e. $d=|g|$ and $\zeta$ is a primitive \emph{$d$-th} root of unity). Now, after Prop.~\ref{Prop: The key technical lemma} again, we must have that $k$ is a multiple of the number $m\coloneqq \frac{d}{\op{gcd}(c_{\chi},d)}$ (since $\Psi^k(\chi)(g)=\zeta^{-kc_{\chi}}\cdot\chi(g)$). Moreover, by Prop.~\ref{Prop: properties of Psi for Cox numbers} the degrees $\chi(1)$ as well as the Coxeter numbers $c_{\chi}$ are not affected by $\Psi$, so that to prove the lemma, it is sufficient to show that if $c_{\chi}$ is not a multiple of $|g|$, then $$\sum_{j=1}^k\Psi^j(\chi)(g^{-1})=0.$$ But if $\xi=\zeta^{-c_{\chi}}$, we have by Prop.~\ref{Prop: The key technical lemma} that $\Psi^j(\chi)(g^{-1})=\xi^j\chi(g^{-1})$ after which the above is immediate (indeed, we have that $m\neq 1$ and $\xi$ is a primitive $m$-th root of unity; therefore $\xi$ is
 also a $k$-th root of unity different from $1$, so that $\sum_{j=1}^k\xi^j=0$).
\end{proof}
